\documentclass[11pt]{article}
\usepackage[margin=0.65in]{geometry}
\usepackage{lmodern}
\usepackage{microtype}
\usepackage{enumitem}
\usepackage{tabularx}
\usepackage[table]{xcolor}
\usepackage{titlesec}
\usepackage[hidelinks]{hyperref}
\setlength{\parindent}{0pt}
\setlength{\parskip}{0.28em}
\setlist[itemize]{leftmargin=1.3em,itemsep=0.08em,topsep=0.1em}
\titleformat{\section}{\large\bfseries}{}{0pt}{}
\titlespacing*{\section}{0pt}{0.65em}{0.2em}
\pagestyle{plain}
\newcommand{\blankline}{\par\vspace{0.12em}\noindent\rule{\linewidth}{0.4pt}\par\vspace{0.38em}}
\newcommand{\smallblank}{\rule{0.43\linewidth}{0.4pt}}
\begin{document}
\begin{center}
{\Large\bfseries U1L10SP --- Short Paper 2}\par
{\LARGE\bfseries Must Caesar Die?}\par
{\large Week 10: Irrelevant Proof and the Point at Issue}
\end{center}
\noindent\textbf{Name:} \smallblank\hfill
\textbf{Date:} \rule{1.35in}{0.4pt}

\begingroup\small
\vspace{0.25em}
\fcolorbox{black}{gray!8}{\begin{minipage}{0.94\linewidth}
\textbf{The question}

In \textit{Julius Caesar}, Brutus concludes, ``It must be by his death.'' Does Brutus prove that Caesar must die, or does he prove only that Caesar might become dangerous? Make and defend your own judgment.

You may conclude that Brutus proves the necessity of killing Caesar, that he proves only a possible danger, or that his case reaches some careful position between those answers. Your score depends on the fit among claim, reasoning, and evidence---not agreement with a preferred answer.
\end{minipage}}

\section*{Text support}
The accompanying \textit{U1L10SP Text Supplement} prints Brutus's orchard reasoning from Act 2, Scene 1 and Antony's counterpressure from Act 3, Scene 2. These are lines worth focusing on, not limits on your evidence. You may use any accurate line, action, choice, or consequence from the play. Quote exact words printed in the supplement; other remembered events may be paraphrased accurately.

\section*{The assignment}
Write a short paper of \textbf{one to two handwritten pages} on separate lined paper. Give it the title \textit{Must Caesar Die?} Your paper should:
\begin{itemize}
\item answer the exact question with a clear, appropriately limited claim;
\item reconstruct Brutus's movement from possible future danger to present death;
\item distinguish what his reasons actually support from what his conclusion requires;
\item use at least two brief quotations or exact details, including evidence from both passages in the supplement;
\item address one reasonable pressure against your judgment and end with a qualification if one is needed.
\end{itemize}
You do not need a five-paragraph essay. Organize the paper in the order that makes your reasoning easiest to follow.

\section*{The closed-world rule}
You may use only the supplied text supplement, the assigned play or course-provided primary text, course handouts, and notes you made during class.

\textbf{Do not use AI, the internet, online summaries, outside criticism, a tutor, a parent, a friend, or anyone else's planning, revising, correcting, sentences, or ideas.} Uncertainty is part of the work; reason through it yourself.

At the end of your paper, write and sign: \emph{I wrote this paper without outside help.}

\section*{What counts --- 10 points}
\rowcolors{2}{gray!8}{white}
\begin{tabularx}{\linewidth}{>{\bfseries}p{0.28\linewidth}Xr}
Claim & Answers the exact question with an appropriately sized judgment. & 2\\
Logical work & Tests the move from possibility to necessity and identifies what remains unproved. & 3\\
Textual evidence & Uses accurate evidence from both supplied passages. & 2\\
Explanation & Shows how each piece of evidence bears on the point at issue. & 2\\
Completion & Is legible, handwritten, and within the assigned length. & 1\\
\end{tabularx}
\endgroup

\newpage
\begin{center}
{\Large\bfseries Planning Scaffold}\par
{\large U1L10SP --- A shorter map for Short Paper 2}
\end{center}
\textbf{Use this page to test your argument, not to fill every space.} Notes and fragments are enough. Your final paper goes on separate lined paper.

\section*{1. Fix the point at issue}
Complete both ideas in your own words.

\textbf{What Brutus must prove:}
\blankline
\textbf{A nearby claim that would not by itself be enough:}
\blankline

\section*{2. Reconstruct and test}
Write Brutus's reasoning as a short chain: what he knows, what he predicts, what he concludes, and what unstated bridge the conclusion needs.
\blankline
\blankline
\blankline

Which link has the strongest evidence? Which link needs more evidence?
\blankline
\blankline

\section*{3. Put the passages under pressure}
Choose one exact line from the orchard passage and one exact line or detail from Antony's speech. Beside each, note the precise claim it supports. Do not ask the quotation to prove more than it can.
\blankline
\blankline
\blankline

\section*{4. Decide the size of your conclusion}
What is the strongest fair objection to your current judgment?
\blankline
What limit word or phrase---such as \emph{might}, \emph{unless}, \emph{not yet}, \emph{only}, or \emph{partly}---would make your conclusion exact?
\blankline

\section*{5. Choose an order}
Number these moves in the order your paper needs: state your claim; reconstruct Brutus; test the gap; weigh the counterpressure; give the qualified judgment.

\vfill
\textbf{Before submitting:} Underline your main claim once. Circle the two places where you use evidence from the play. Sign the authorship statement.
\end{document}
